\chapter{Fossil-CMB Thermalization and Active-Vacuum Equilibrium}
\label{app:cmb_thermal_completion}

\noindent\textbf{Abstract.}
This appendix examines the thermodynamic completion required to connect the
coherent photon scale-transfer mechanism developed in
\hyperref[app:photon_dispersion_relaxation]{Appendix~\ref*{app:photon_dispersion_relaxation}}
and the microscopic relaxation-rate construction developed in
\hyperref[app:vacuumdamping_cosmicredshift]{Appendix~\ref*{app:vacuumdamping_cosmicredshift}}
to a physically consistent Cosmic Microwave Background (CMB) sector.

The analysis separates four mechanically distinct processes:
coherent photon scale transfer, excitation of noncoherent finite-energy
Active-Vacuum degrees of freedom, thermalization with photon-number-changing
detailed balance, and macroscopic thermal energy balance. The inherited
coherent law
\(
\dot E_\gamma=-\Gamma_{\mathrm D}E_\gamma
\)
may be represented as the coarse-grained first moment of complete quantized
energy-transfer events and therefore need not imply fractional quantum
occupation of individual receiving modes.

Appendix~\ref{app:vacuumdamping_cosmicredshift} develops a conditional
microscopic construction of the relaxation rate. A nonzero constant
\(\Gamma_{\mathrm D}\) requires an independently established
carrier-projected irreversible continuum rather than a finite conservative
pure-point mode set. The present appendix addresses the separate
thermodynamic question of what density of states and heat capacity follow if
finite-energy noncoherent receiving branches exist.

A central distinction is established between the thermodynamic density of
states and the carrier-projected spectral measure governing irreversible
photon coupling. Even complete knowledge of the latter does not determine
the thermodynamic multiplicity or heat capacity of the receiver. For
candidate three-dimensional gapless linear bosonic branches, the
low-temperature energy density scales as \(T^4\) and the corresponding heat
capacity scales as \(T^3\).

A separate photon-number-changing absorption--emission interaction satisfying
microscopic detailed balance is required for a zero-chemical-potential Planck
occupation. The thermodynamic receiving spectrum, microscopic
photon--receiver coupling, primordial receiver loading, thermal interaction,
and absolute temperature normalization must therefore be derived
independently rather than calibrated from the observed CMB temperature.


\section{Purpose and Inherited Appendix Hierarchy}

The purpose of this appendix is to determine what additional thermodynamic
structure is required once coherent photon energy transfer has been
specified.

The inherited coherent photon law is
\begin{equation}
    \frac{\mathrm dE_\gamma}{\mathrm dt}
    =
    -\Gamma_{\mathrm D}E_\gamma,
    \qquad
    \frac{\mathrm d\omega_\gamma}{\mathrm dt}
    =
    -\Gamma_{\mathrm D}\omega_\gamma.
    \label{eq:AT_coherent_transport}
\end{equation}

For a quantum-normalized surviving photon,
\begin{equation}
    I_\gamma
    =
    \frac{E_\gamma}{\omega_\gamma}
    =
    \hbar.
    \label{eq:AT_action}
\end{equation}

A conditional microscopic construction of
\(\Gamma_{\mathrm D}\)
is developed in
Appendix~\ref{app:vacuumdamping_cosmicredshift}.
A nonzero constant rate requires an independently established
carrier-projected irreversible continuum. In the absence of such a
continuum, the completed conservative microscopic parent model does not
generate irreversible damping of the pristine coherent photon branch.

The present appendix neither recalibrates
\(\Gamma_{\mathrm D}\)
nor identifies its carrier-projected dissipative response with a
thermodynamic density of states.

The propagation contribution to astronomical redshift is
\(z_{\mathrm D}\), not the complete observed redshift. For stationary
gravitational conditions and consistently factorized endpoint comparisons,
Appendix~\ref{app:time_dilation} gives
\begin{equation}
    \boxed{
    1+z_{\mathrm{obs}}
    =
    (1+z_{\mathrm{kin}})
    (1+z_{\mathrm{grav}})
    (1+z_{\mathrm D}).
    }
    \label{eq:AT_redshift_hierarchy}
\end{equation}

In the gravity-neutral endpoint limit,
\begin{equation}
    z_{\mathrm{grav}}
    \rightarrow
    0,
\end{equation}
Equation~\eqref{eq:AT_redshift_hierarchy} reduces to the corresponding
two-factor kinematic--relaxation form.

For a positive constant physical relaxation rate, the characteristic
coherent-relaxation scale is
\begin{equation}
    R_{\mathrm U}
    =
    \frac{c_{\mathrm E}}{\Gamma_{\mathrm D}},
    \label{eq:AT_RU}
\end{equation}
with
\begin{equation}
    R_{\mathrm U}
    \neq
    R_{\mathrm{Res}}.
    \label{eq:AT_RU_RRes}
\end{equation}

The Hard Floor remains the zero-inertia constraint, phase-reference, and
information-memory sector. It is not assigned ordinary finite thermal
inertia and is not used as a thermal reservoir in the present construction.

The logical hierarchy required for CMB completion is therefore
\begin{equation}
    \begin{aligned}
        \text{coherent photon}
        &\longrightarrow
        \text{noncoherent finite-energy receiving modes}
        \\
        &\longrightarrow
        \text{thermalized Active-Vacuum sector}
        \\
        &\rightleftharpoons
        \text{thermal photon field}
        \\
        &\longrightarrow
        \text{macroscopic thermal balance}.
    \end{aligned}
    \label{eq:AT_hierarchy}
\end{equation}

Each arrow denotes a distinct mechanical requirement.

In particular, coherent energy transfer does not by itself establish
thermalization; thermalization does not by itself enforce a Planck spectrum;
and the existence of a Planck spectrum does not by itself determine an
absolute temperature.


\section{Quantum-Consistent Receiving Channel}

The differential law in
Equation~\eqref{eq:AT_coherent_transport}
must not be interpreted as depositing an infinitesimal fraction of
\(\hbar\)
on an individual microscopic degree of freedom.

Let
\(\Omega_\lambda\)
denote the frequencies of noncoherent finite-energy receiving modes.
A forward quantum-consistent transfer event may be represented schematically
as
\begin{equation}
    \left|
        1_{\mathbf k};
        \{n_\lambda\}
    \right\rangle
    \longrightarrow
    \left|
        1_{\mathbf k'};
        \{n_\lambda+\Delta n_\lambda\}
    \right\rangle,
    \label{eq:AT_quantized_transition}
\end{equation}
with
\begin{equation}
    \Delta n_\lambda
    =
    0,1,2,\ldots
\end{equation}
for the forward energy-transfer process.

Energy conservation requires
\begin{equation}
    \hbar\omega_{\mathbf k}
    -
    \hbar\omega_{\mathbf k'}
    =
    \sum_\lambda
    \Delta n_\lambda
    \hbar\Omega_\lambda.
    \label{eq:AT_quantized_energy}
\end{equation}

The selected coherent carrier survives the event. Its photon number remains
unity,
\begin{equation}
    N_{\mathrm{ph}}
    =
    1,
\end{equation}
while
\begin{equation}
    \frac{E_\gamma}{\omega_\gamma}
    =
    \hbar
\end{equation}
is retained.

The notation
\(N_{\mathrm{ph}}\)
is used specifically for photon number and is distinct from information-count
quantities defined elsewhere in TFT.

Let
\(W(\Omega|\omega)\,\mathrm d\Omega\)
denote the transition rate associated with transferred frequencies between
\(\Omega\)
and
\(\Omega+\mathrm d\Omega\).
The mean coherent-energy drift is
\begin{equation}
    -\frac{\mathrm dE_\gamma}{\mathrm dt}
    =
    \int_0^\omega
    \hbar\Omega
    W(\Omega|\omega)
    \,\mathrm d\Omega.
    \label{eq:AT_first_jump_moment}
\end{equation}

Consistency with
Equation~\eqref{eq:AT_coherent_transport}
requires
\begin{equation}
    \int_0^\omega
    \Omega
    W(\Omega|\omega)
    \,\mathrm d\Omega
    =
    \Gamma_{\mathrm D}\omega.
    \label{eq:AT_drift_constraint}
\end{equation}

The second jump-moment rate is
\begin{equation}
    \mathcal B_{\mathrm E}^{(2)}(\omega)
    =
    \int_0^\omega
    (\hbar\Omega)^2
    W(\Omega|\omega)
    \,\mathrm d\Omega.
    \label{eq:AT_second_jump_moment}
\end{equation}

Higher jump moments are defined analogously.

A viable microscopic realization must therefore reproduce the first moment
in
Equation~\eqref{eq:AT_drift_constraint}
while keeping the accumulated second and higher moments compatible with the
spectral-coherence constraints developed in
Appendices~\ref{app:photon_dispersion_relaxation}
and~\ref{app:time_dilation}.

A sparse sequence of comparatively large inelastic events would generically
produce appreciable stochastic frequency diffusion. The required coherent
limit is instead one in which the mean drift remains finite while the
associated spectral variance remains sufficiently small.


\section{Thermodynamic Receiving Spectrum and Carrier Projection}

Appendix~\ref{app:vacuumdamping_cosmicredshift}
establishes the dynamical requirement that a nonzero irreversible
low-frequency photon response cannot be generated by a finite conservative
pure-point mode set alone.

The present section asks instead what thermodynamic state counting follows if
a specified finite-energy receiving continuum exists.

Consider a candidate receiving branch whose low-wave-number dispersion obeys
\begin{equation}
    \Omega_j(\mathbf k_{\mathrm{rec}})
    \propto
    k_{\mathrm{rec}}^{p_j},
    \qquad
    p_j>0,
    \label{eq:AT_general_dispersion}
\end{equation}
where
\begin{equation}
    k_{\mathrm{rec}}
    =
    \left|
        \mathbf k_{\mathrm{rec}}
    \right|
\end{equation}
is the dimensional wave number of the receiving branch. It is distinct from
the dimensionless photon scale coordinate
\(q=kr_{\mathrm D}\)
used in Appendix~\ref{app:vacuumdamping_cosmicredshift}.

Let
\(\mathcal D_j(\Omega)\)
denote its thermodynamic density of states.

In three physical spatial dimensions,
\begin{equation}
    \frac{\mathcal D_j(\Omega)}{V}
    \propto
    \Omega^{3/p_j-1}.
    \label{eq:AT_general_DOS}
\end{equation}

For a gapless linear branch,
\begin{equation}
    \Omega_j(\mathbf k_{\mathrm{rec}})
    =
    v_jk_{\mathrm{rec}},
    \label{eq:AT_linear_dispersion}
\end{equation}
the density of states per physical volume is
\begin{equation}
    \frac{\mathcal D_j(\Omega)}{V}
    =
    \frac{\Omega^2}{2\pi^2v_j^3}.
    \label{eq:AT_linear_DOS}
\end{equation}

For several independent branches,
\begin{equation}
    \frac{\mathcal D_{\mathrm{AV}}(\Omega)}{V}
    =
    \frac{\Omega^2}{2\pi^2}
    \sum_j
    v_j^{-3}.
    \label{eq:AT_total_DOS}
\end{equation}

Equation~\eqref{eq:AT_total_DOS}
is a state-counting result.

It does not determine the strength with which a coherent photon couples to
those states.

Define
\begin{equation}
    \mathcal S_{\gamma\mathrm{AV}}(\Omega)
\end{equation}
as the carrier-projected Active-Vacuum spectral measure entering the
irreversible photon response.

The thermodynamic and dissipative quantities are therefore distinct:
\begin{equation}
    \mathcal D_{\mathrm{AV}}(\Omega)
    \not\equiv
    \mathcal S_{\gamma\mathrm{AV}}(\Omega).
    \label{eq:AT_DOS_spectral_distinction}
\end{equation}

If the projected spectral measure admits a factorized representation of the
schematic form
\begin{equation}
    \mathcal S_{\gamma\mathrm{AV}}(\Omega)
    \propto
    \frac{\mathcal D_{\mathrm{AV}}(\Omega)}{V}
    \left|
        \langle f|
        H_{\mathrm{int}}
        |i\rangle
    \right|^2,
    \label{eq:AT_projected_spectrum}
\end{equation}
then both the density of states and the low-frequency interaction matrix
element determine the irreversible response.

For a frequency-independent low-frequency damping coefficient, the
carrier-projected response must satisfy the Ohmic scaling
\begin{equation}
    \mathcal S_{\gamma\mathrm{AV}}(\Omega)
    \propto
    \Omega,
    \qquad
    \Omega\rightarrow0.
    \label{eq:AT_ohmic_requirement}
\end{equation}

For the three-dimensional linear density of states in
Equation~\eqref{eq:AT_linear_DOS},
the factorized representation would therefore require
\begin{equation}
    \left|
        \langle f|
        H_{\mathrm{int}}
        |i\rangle
    \right|^2
    \propto
    \Omega^{-1}
    \label{eq:AT_linear_vertex_requirement}
\end{equation}
in the relevant low-frequency limit.

More generally, for
\(\Omega_j\propto k_{\mathrm{rec}}^{p_j}\),
\begin{equation}
    \left|
        \langle f|
        H_{\mathrm{int}}
        |i\rangle
    \right|^2
    \propto
    \Omega^{\,2-3/p_j}
    \label{eq:AT_general_vertex_requirement}
\end{equation}
would be required within the same factorized representation.

Equations~\eqref{eq:AT_linear_vertex_requirement}
and~\eqref{eq:AT_general_vertex_requirement}
are scaling requirements on a hypothetical factorized interaction.
They do not derive the microscopic interaction itself.

Accordingly,
\begin{equation}
    \text{thermal state counting}
    \not\Rightarrow
    \text{irreversible photon coupling}.
    \label{eq:AT_state_count_not_coupling}
\end{equation}


\subsection{Thermodynamic Multiplicity Is Not Fixed by Carrier Projection}

The distinction between thermodynamic state counting and carrier-projected
dissipation remains even if
\(\mathcal S_{\gamma\mathrm{AV}}\)
is known.

Within the factorized representation of
Equation~\eqref{eq:AT_projected_spectrum},
consider the simultaneous rescaling
\begin{equation}
    \mathcal D_{\mathrm{AV}}(\Omega)
    \longrightarrow
    \mathcal N_{\mathrm{mult}}
    \mathcal D_{\mathrm{AV}}(\Omega),
\end{equation}
and
\begin{equation}
    \left|
        \langle f|
        H_{\mathrm{int}}
        |i\rangle
    \right|^2
    \longrightarrow
    \frac{1}{\mathcal N_{\mathrm{mult}}}
    \left|
        \langle f|
        H_{\mathrm{int}}
        |i\rangle
    \right|^2,
\end{equation}
where
\(\mathcal N_{\mathrm{mult}}>0\).

The carrier-projected spectral measure is unchanged, while the thermal energy
associated with the receiving density of states is multiplied by
\(\mathcal N_{\mathrm{mult}}\).

Therefore,
\begin{equation}
    \mathcal S_{\gamma\mathrm{AV}}(\Omega)
    \not\Rightarrow
    u_{\mathrm{AV}}(T).
    \label{eq:AT_spectrum_not_thermal_energy}
\end{equation}

Thus a dissipative spectral measure does not determine receiver
thermodynamic multiplicity or heat capacity.

The normalized overlap factor
\(P_{\mathrm R}\)
introduced in
Appendix~\ref{app:vacuumdamping_cosmicredshift}
must be interpreted accordingly.

For a genuine projector
\(\Pi_{\mathrm R}\),
\begin{equation}
    \Pi_{\mathrm R}^2
    =
    \Pi_{\mathrm R}.
\end{equation}

If
\begin{equation}
    \widetilde{\Pi}_{\mathrm R}
    =
    P_{\mathrm R}\Pi_{\mathrm R},
\end{equation}
then
\begin{equation}
    \widetilde{\Pi}_{\mathrm R}^{\,2}
    =
    P_{\mathrm R}
    \widetilde{\Pi}_{\mathrm R},
\end{equation}
which is not idempotent unless
\(P_{\mathrm R}=0\)
or
\(P_{\mathrm R}=1\).

Therefore
\(P_{\mathrm R}\)
may enter a coupling or kernel weight, but it cannot by itself represent
thermodynamic projector rank, state count, or degeneracy.

Thermal-temperature estimates obtained by identifying
\(P_{\mathrm R}\)
directly with a thermodynamic mode fraction are not retained.


\section{Conditional Thermal Energy and Heat Capacity}

Suppose a specified set of noncoherent Active-Vacuum receiving modes
thermalizes and obeys the Bose occupation
\begin{equation}
    n_{\mathrm B}(\Omega,T)
    =
    \frac{1}{
        \exp
        \left(
            \hbar\Omega/k_{\mathrm B}T
        \right)
        -1
    }.
    \label{eq:AT_Bose_occupation}
\end{equation}

Their thermal energy density, excluding the temperature-independent
zero-point contribution, is
\begin{equation}
    u_{\mathrm{AV}}(T)
    =
    \frac{1}{V}
    \sum_j
    \int_0^\infty
    \hbar\Omega\,
    \mathcal D_j(\Omega)\,
    n_{\mathrm B}(\Omega,T)
    \,\mathrm d\Omega.
    \label{eq:AT_thermal_energy_general}
\end{equation}

For gapless linear branches,
substitution of
Equation~\eqref{eq:AT_linear_DOS}
gives
\begin{equation}
    u_{\mathrm{AV}}(T)
    =
    \frac{\pi^2}{30}
    \frac{k_{\mathrm B}^4T^4}{\hbar^3}
    \sum_j
    v_j^{-3}.
    \label{eq:AT_thermal_energy_linear}
\end{equation}

Equation~\eqref{eq:AT_thermal_energy_linear}
is the low-temperature continuum result for the selected linear branches.
It applies only while the thermally occupied modes remain in the regime for
which
\(\Omega_j\simeq v_jk_{\mathrm{rec}}\).

If significant dispersion curvature or a finite band edge is sampled, the
exact density-of-states integral in
Equation~\eqref{eq:AT_thermal_energy_general}
must be used instead.

The corresponding volumetric heat capacity in the linear regime is
\begin{equation}
    C_V^{(\mathrm{AV})}(T)
    \equiv
    \frac{\mathrm du_{\mathrm{AV}}}{\mathrm dT}
    =
    \frac{2\pi^2}{15}
    \frac{k_{\mathrm B}^4T^3}{\hbar^3}
    \sum_j
    v_j^{-3}.
    \label{eq:AT_heat_capacity_linear}
\end{equation}

For the more general low-wave-number dispersion
\(\Omega_j\propto k_{\mathrm{rec}}^{p_j}\),
the corresponding low-temperature continuum scalings are
\begin{equation}
    u(T)
    \propto
    T^{1+3/p_j},
    \qquad
    C_V(T)
    \propto
    T^{3/p_j}.
    \label{eq:AT_general_thermal_scaling}
\end{equation}

These are conditional statistical-mechanical results.

They determine the heat capacity only after the physical receiving branches,
their state-counting measure, and their dispersions have been independently
specified.

They do not determine the photon-projected spectral measure
\(\mathcal S_{\gamma\mathrm{AV}}\).

The thermal energy stored in nonradiative receiving modes must also be
distinguished from the electromagnetic CMB energy density:
\begin{equation}
    u_{\mathrm{CMB}}
    \not\equiv
    u_{\mathrm{AV}}
    \label{eq:AT_radiative_thermal_distinction}
\end{equation}
in general.

Here
\(u_{\mathrm{CMB}}\)
denotes the energy density carried by the thermal electromagnetic field,
whereas
\(u_{\mathrm{AV}}\)
denotes thermal energy stored in the specified noncoherent Active-Vacuum
receiving modes.

The total thermalized sector may therefore be written schematically as
\begin{equation}
    u_{\mathrm{th}}
    =
    u_{\mathrm{CMB}}
    +
    u_{\mathrm{nonrad}},
    \label{eq:AT_total_thermal_partition}
\end{equation}
where
\(u_{\mathrm{nonrad}}\)
collects independently established thermal non-electromagnetic degrees of
freedom.

No equality between
\(u_{\mathrm{CMB}}\)
and the complete thermal Active-Vacuum energy density is assumed.

Accordingly, a candidate three-dimensional gapless linear reservoir
establishes conditionally
\begin{equation}
    u_{\mathrm{AV}}\propto T^4,
    \qquad
    C_V^{(\mathrm{AV})}\propto T^3,
\end{equation}
but does not by itself establish the microscopic photon-transfer interaction,
number-changing Planck equilibration, primordial receiver loading, or
absolute CMB temperature.


\section{Continuous Transfer and the Thermal Energy Ledger}

Appendix~\ref{app:photon_dispersion_relaxation}
develops the coherent-carrier propagation law and its associated continuous
and terminal energy bookkeeping.

For the thermodynamic analysis, the required inherited identity is
\begin{equation}
    E_{\mathrm{emit}}
    =
    \Delta E_{\mathrm{cont}}
    +
    \Delta E_{\mathrm{term}},
    \label{eq:AT_total_energy_ledger}
\end{equation}
where
\(\Delta E_{\mathrm{cont}}\)
is the energy transferred continuously from the coherent carrier before any
terminal event and
\(\Delta E_{\mathrm{term}}\)
is the remaining coherent contribution transferred at terminal coherence
loss, if such a terminal event exists.

For a terminal coherent energy
\(E_{\mathrm{term}}\),
\begin{equation}
    \Delta E_{\mathrm{cont}}
    =
    E_{\mathrm{emit}}
    -
    E_{\mathrm{term}},
    \qquad
    \Delta E_{\mathrm{term}}
    =
    E_{\mathrm{term}}.
    \label{eq:AT_energy_partition}
\end{equation}

Hence
\begin{equation}
    E_{\mathrm{emit}}
    \gg
    E_{\mathrm{term}}
    \quad
    \Longrightarrow
    \quad
    \Delta E_{\mathrm{cont}}
    \gg
    \Delta E_{\mathrm{term}}.
    \label{eq:AT_terminal_subdominant_limit}
\end{equation}

A terminal coherence-loss event therefore cannot generically be treated as
the sole destination of the initial photon energy when most of that energy
has already left the coherent carrier continuously.

The continuously transferred contribution must remain in the complete energy
ledger. It may populate nonthermal receiving excitations, thermalized
Active-Vacuum modes, or another independently derived finite-energy sector.

Accordingly,
\begin{equation}
    \text{coherent energy transfer}
    \not\Rightarrow
    \text{immediate thermalization}.
    \label{eq:AT_transfer_not_thermalization}
\end{equation}

The microscopic terminal condition, if one exists, must remain independent
of the observed CMB spectrum.

The required logical ordering is
\begin{equation}
    \text{photon-branch stability}
    \longrightarrow
    E_{\mathrm{term}}^{(\mathrm{TFT})}
    \longrightarrow
    \text{thermal evolution}
    \longrightarrow
    \text{CMB comparison}.
    \label{eq:AT_terminal_logic}
\end{equation}

Neither
\(E_{\mathrm{term}}\)
nor any corresponding terminal frequency is used below as an upstream
normalization of the CMB temperature.


\section{Why Coherent Scale Transfer Alone Does Not Produce a Blackbody}

The coherent scale-transfer branch preserves the identity of the surviving
photon.

Its photon-number condition is therefore
\begin{equation}
    \dot N_{\mathrm{ph}}
    =
    0
    \label{eq:AT_coherent_number}
\end{equation}
for that branch.

This condition is distinct from the occupation-number evolution of a
subsequently thermalized electromagnetic population.

A bosonic population that reaches kinetic equilibrium while its particle
number remains independently constrained is described generally by the
Bose--Einstein occupation
\begin{equation}
    n_{\mathrm{BE}}(\omega,T)
    =
    \frac{1}{
        \exp
        \left[
            \left(
                \hbar\omega-\mu_\gamma
            \right)
            /
            k_{\mathrm B}T
        \right]
        -1
    },
    \label{eq:AT_BE_mu}
\end{equation}
where
\(\mu_\gamma\)
is the photon chemical potential associated with the number constraint.

The Planck occupation is the special case
\begin{equation}
    \mu_\gamma
    =
    0.
    \label{eq:AT_zero_mu}
\end{equation}

Consequently,
\begin{equation}
    \text{coherent frequency redistribution with fixed }
    N_{\mathrm{ph}}
    \not\Rightarrow
    \text{Planck equilibrium}.
    \label{eq:AT_number_conservation_no_Planck}
\end{equation}

This conclusion is independent of the detailed microscopic origin of
\(\Gamma_{\mathrm D}\).

A separate thermal interaction must permit the electromagnetic population to
change occupation number so that
\(\mu_\gamma\)
can relax to zero.

The coherent and thermal sectors therefore obey different number
constraints:
\begin{equation}
    \begin{aligned}
        \text{coherent branch:}
        &\qquad
        \dot N_{\mathrm{ph}}=0,
        \\
        \text{thermal branch:}
        &\qquad
        \dot N_{\mathrm{ph}}
        \text{ need not vanish}.
    \end{aligned}
    \label{eq:AT_number_sector_distinction}
\end{equation}


\section{Number-Changing Detailed Balance and the Conditional Planck State}

Let
\(n(\omega,t)\)
denote the occupation number of a thermal electromagnetic mode.

A minimal effective absorption--emission collision operator may be written
as
\begin{equation}
    C_{\mathrm{th}}[n]
    =
    \Gamma_{\mathrm{em}}(\omega,T)
    \left(
        1+n
    \right)
    -
    \Gamma_{\mathrm{abs}}(\omega,T)
    n.
    \label{eq:AT_collision_operator}
\end{equation}

Here
\(\Gamma_{\mathrm{em}}\)
is the effective emission rate into the electromagnetic mode and
\(\Gamma_{\mathrm{abs}}\)
is the corresponding absorption rate.

Microscopic detailed balance at temperature \(T\) requires
\begin{equation}
    \frac{
        \Gamma_{\mathrm{em}}(\omega,T)
    }{
        \Gamma_{\mathrm{abs}}(\omega,T)
    }
    =
    \exp
    \left(
        -
        \frac{
            \hbar\omega
        }{
            k_{\mathrm B}T
        }
    \right).
    \label{eq:AT_detailed_balance}
\end{equation}

At equilibrium,
\begin{equation}
    C_{\mathrm{th}}[n_*]
    =
    0.
    \label{eq:AT_collision_equilibrium}
\end{equation}

Substitution of
Equation~\eqref{eq:AT_detailed_balance}
into
Equation~\eqref{eq:AT_collision_operator}
gives
\begin{equation}
    n_*(\omega,T)
    =
    \frac{1}{
        \exp
        \left(
            \hbar\omega/k_{\mathrm B}T
        \right)
        -1
    }.
    \label{eq:AT_Planck_occupation}
\end{equation}

Thus a zero-chemical-potential Planck occupation follows conditionally if the
thermalized sector possesses a genuine photon-number-changing
absorption--emission interaction satisfying microscopic detailed balance.

At the microscopic level, that interaction must satisfy
\begin{equation}
    \left[
        \hat H_{\mathrm{th}},
        \hat N_{\mathrm{ph}}
    \right]
    \neq
    0.
    \label{eq:AT_number_changing_requirement}
\end{equation}

This interaction is physically distinct from the coherent scale-transfer
interaction, which preserves the identity of the surviving photon.

Linearizing
Equation~\eqref{eq:AT_collision_operator}
about
\begin{equation}
    n
    =
    n_{\mathrm P}
    +
    \delta n
\end{equation}
gives
\begin{equation}
    \frac{\mathrm d\delta n}{\mathrm dt}
    =
    -
    \left(
        \Gamma_{\mathrm{abs}}
        -
        \Gamma_{\mathrm{em}}
    \right)
    \delta n.
    \label{eq:AT_Planck_stability}
\end{equation}

For finite positive
\(\omega/T\),
Equation~\eqref{eq:AT_detailed_balance}
implies
\begin{equation}
    \Gamma_{\mathrm{abs}}
    >
    \Gamma_{\mathrm{em}},
\end{equation}
so the Planck solution is linearly stable within this effective kinetic
description.

Define
\begin{equation}
    \tau_{\mathrm{th}}^{-1}
    =
    \Gamma_{\mathrm{abs}}
    -
    \Gamma_{\mathrm{em}}.
    \label{eq:AT_thermalization_rate}
\end{equation}

A sufficient hierarchy for rapid restoration of the thermal spectrum
relative to coherent cosmological scale transfer is
\begin{equation}
    \tau_{\mathrm{th}}
    \ll
    \Gamma_{\mathrm D}^{-1}.
    \label{eq:AT_thermalization_hierarchy}
\end{equation}

This inequality is a consistency condition and does not determine
\(\tau_{\mathrm{th}}\).

The thermal number-changing interaction must simultaneously remain
sufficiently selective that photons remaining on the coherent astronomical
branch do not acquire unacceptable absorption, scattering, spectral
broadening, or polarization conversion.

That selectivity remains a microscopic completion requirement.

A strict equilibrium collision operator satisfies
\begin{equation}
    C_{\mathrm{th}}[n_{\mathrm P}]
    =
    0.
    \label{eq:AT_equilibrium_collision_zero}
\end{equation}

Therefore the coherent scale-transfer law and the equilibrium thermal
collision operator should not be applied as two independent evolution laws
to the same mode population unless a driven nonequilibrium kinetic equation
has been explicitly derived.

The retained sector separation is
\begin{equation}
    \text{coherent photon branch}
    \qquad\longrightarrow\qquad
    \text{thermal photon sector}.
    \label{eq:AT_sector_separation}
\end{equation}


\section{Fossil Initial-State Boundary}

The Fossil-CMB route treats the present thermal radiation field as a relic of
the Catastrophic Decompression and annealing history rather than as radiation
generated by a present-day steady-state source.

The high-temperature quantity
$T_{\mathrm S}$
introduced in
\hyperref[app:lattice_annealing]{Appendix~\ref*{app:lattice_annealing}}
is the kinetic-equivalent temperature scale associated with the substrate
phase transition and is not identified directly with an ordinary photon
blackbody temperature.

The earliest presently defensible thermal photon boundary is the
annealing/transparency endpoint derived in
\hyperref[app:lattice_annealing]{Appendix~\ref*{app:lattice_annealing}},
\begin{equation}
    t_{\mathrm a}
    \simeq
    3.86\times10^5~{\rm yr},
    \qquad
    T_{\gamma}(t_{\mathrm a})
    =
    T_{\mathrm a}
    \sim
    3000~{\rm K}.
    \label{eq:AT_transparency_temperature}
\end{equation}

The same pair $(T_{\mathrm a},t_{\mathrm a})$ therefore defines the endpoint
of lattice annealing and the initial boundary of the subsequent Fossil-CMB
thermal evolution.

A Planck spectrum at this boundary remains conditional on sufficiently rapid
early photon-number-changing thermalization.

The high-temperature annealing similarity relation used in
\hyperref[app:lattice_annealing]{Appendix~\ref*{app:lattice_annealing}}
is retained within the regime for which it was derived and is not extrapolated
as a post-transparency CMB cooling law.

In particular,
\begin{equation}
    T\propto t^{-1/2}
\end{equation}
does not by itself determine the subsequent temperature of a finite fossil
radiation field.

Likewise, long propagation distance or random spatial wandering does not
constitute a thermal cooling law.

For a homogeneous thermal radiation field in a fixed physical volume,
\begin{equation}
    \frac{\mathrm d u_{\mathrm{CMB}}}{\mathrm dt}
    =
    \mathcal P_{\mathrm{dep}}
    -
    \mathcal P_{\mathrm{loss}},
    \label{eq:AT_fossil_first_law}
\end{equation}
where
$\mathcal P_{\mathrm{dep}}$
and
$\mathcal P_{\mathrm{loss}}$
are respectively the net power densities deposited into and removed from the
thermal photon sector.

Since
\begin{equation}
    u_{\mathrm{CMB}}\propto T^4,
\end{equation}
a system satisfying
\begin{equation}
    \mathcal P_{\mathrm{loss}}=0,
    \qquad
    \mathcal P_{\mathrm{dep}}\geq0
\end{equation}
cannot cool:
\begin{equation}
    \dot T\geq0.
    \label{eq:AT_no_loss_no_cooling}
\end{equation}

Thus post-transparency Fossil cooling requires a genuine finite-energy
receiving channel or another explicitly derived energy-loss mechanism.
Reduction of active FCC volume is likewise not identified with ordinary
adiabatic expansion cooling.

\section{Conditional Fossil-Radiation Energy-Conservation Closure}

The preceding sections establish conditions under which a thermal
electromagnetic population may possess a Planck spectrum.

They do not determine the temperature of that spectrum.

Consider a fixed physical coarse-graining volume during an equilibration
interval in which external work, macroscopic volume change, and energy
exchange with any third sector are negligible.

Let
\(T_{\gamma,0}\)
denote the initial photon temperature and
\(T_{\mathrm{rec},0}\)
the initial effective temperature of the selected receiving sector.

The subscript ``0'' labels the initial state of this equilibration problem and
does not denote the TFT Hard-Floor State~0.

If complete mutual thermal equilibrium is reached, let the common final
temperature be
\(T_{\mathrm{eq}}\).

For an isotropic electromagnetic field with two transverse polarizations and
propagation speed
\(c_{\mathrm E}\),
\begin{equation}
    u_{\mathrm{CMB}}(T)
    =
    \frac{
        \pi^2k_{\mathrm B}^4
    }{
        15\hbar^3c_{\mathrm E}^3
    }
    T^4.
    \label{eq:AT_CMB_energy_density}
\end{equation}

For independent three-dimensional gapless linear receiving branches with
speeds
\(v_j\),
Equation~\eqref{eq:AT_thermal_energy_linear}
gives
\begin{equation}
    u_{\mathrm{rec}}(T)
    =
    \frac{
        \pi^2k_{\mathrm B}^4
    }{
        30\hbar^3
    }
    \left(
        \sum_j v_j^{-3}
    \right)
    T^4.
    \label{eq:AT_receiver_energy_density}
\end{equation}

This expression applies only while the thermally occupied portion of each
branch remains within its approximately linear continuum regime.

For an isolated photon--receiver subsystem, energy conservation requires
\begin{equation}
    u_{\mathrm{CMB}}
    \left(
        T_{\gamma,0}
    \right)
    +
    u_{\mathrm{rec}}
    \left(
        T_{\mathrm{rec},0}
    \right)
    =
    u_{\mathrm{CMB}}
    \left(
        T_{\mathrm{eq}}
    \right)
    +
    u_{\mathrm{rec}}
    \left(
        T_{\mathrm{eq}}
    \right).
    \label{eq:AT_fossil_energy_conservation_general}
\end{equation}

Substitution gives
\begin{equation}
    T_{\mathrm{eq}}^4
    =
    \frac{
        2c_{\mathrm E}^{-3}
        T_{\gamma,0}^4
        +
        \left(
            \displaystyle\sum_j v_j^{-3}
        \right)
        T_{\mathrm{rec},0}^4
    }{
        2c_{\mathrm E}^{-3}
        +
        \displaystyle\sum_j v_j^{-3}
    }.
    \label{eq:AT_fossil_temperature_general}
\end{equation}

Equation~\eqref{eq:AT_fossil_temperature_general}
is a conditional consequence of energy conservation and the specified
receiving thermodynamics.

It does not determine the receiving branches, their propagation speeds, their
initial loading, or the microscopic rate of photon--receiver exchange.

For one distinct gapless scalar receiving branch with speed
\(v_{\mathrm{rec}}\),
\begin{equation}
    T_{\mathrm{eq}}^4
    =
    \frac{
        T_{\gamma,0}^4
        +
        \dfrac{1}{2}
        \left(
            \dfrac{c_{\mathrm E}}{v_{\mathrm{rec}}}
        \right)^3
        T_{\mathrm{rec},0}^4
    }{
        1
        +
        \dfrac{1}{2}
        \left(
            \dfrac{c_{\mathrm E}}{v_{\mathrm{rec}}}
        \right)^3
    }.
    \label{eq:AT_fossil_temperature_scalar}
\end{equation}

The notation
\(v_{\mathrm{rec}}\)
is intentionally generic at this stage.

No identification of the general receiver with an ordinary FCC phonon, the
Hard Floor, a bulk FCC--SC conversion coordinate, or a particular internal
phase branch is implied by
Equation~\eqref{eq:AT_fossil_temperature_scalar}.


\subsection{Consistency Limits and Receiver-State Separation}

If
\begin{equation}
    T_{\mathrm{rec},0}
    =
    T_{\gamma,0},
\end{equation}
then
\begin{equation}
    T_{\mathrm{eq}}
    =
    T_{\gamma,0}.
    \label{eq:AT_equal_temperature_check}
\end{equation}

Thus a receiving sector already equilibrated with the thermal radiation at the
initial boundary cannot subsequently cool that radiation through closed mutual
equilibration.

The modes strongly populated during catastrophic decompression and early
annealing must therefore be distinguished from any weakly populated receiver
responsible for later Fossil evolution:
\begin{equation}
    \mathcal H_{\mathrm{AP}}
    \neq
    \mathcal H_{\mathrm{rec}}.
    \label{eq:AT_AP_receiver_separation}
\end{equation}

This separation is a microscopic requirement rather than a thermodynamic
identity.

In the limiting cold-receiver case,
\begin{equation}
    T_{\mathrm{rec},0}
    =
    0,
\end{equation}
Equation~\eqref{eq:AT_fossil_temperature_scalar}
becomes
\begin{equation}
    T_{\mathrm{eq}}
    =
    T_{\gamma,0}
    \left[
        1+
        \frac{1}{2}
        \left(
            \frac{c_{\mathrm E}}{v_{\mathrm{rec}}}
        \right)^3
    \right]^{-1/4}.
    \label{eq:AT_cold_receiver_limit}
\end{equation}

This is a boundary-case solution and does not establish that the primordial
receiver was actually cold.


\subsection{Primordial Receiver Loading}

Weak photon coupling before transparency does not by itself imply a weakly
populated receiving sector.

Using the unit-reference relaxation rate
\(\Gamma_{\mathrm D}^{(0)}\)
together with the retained annealing/transparency age gives the illustrative
pre-transparency exposure
\begin{equation}
    \Gamma_{\mathrm D}^{(0)}t_{\mathrm a}
    \simeq
    2.87\times10^{-5}.
    \label{eq:AT_pretransparency_exposure}
\end{equation}

This quantity constrains only loading acquired through the corresponding
reference photon-coupling channel.

Catastrophic decompression may excite the receiving sector directly. Its
microscopic energy ledger therefore has the schematic form
\begin{equation}
    \dot u_{\mathrm{rec}}
    =
    \mathcal P_{\mathrm{form,rec}}
    +
    \mathcal P_{\gamma\mathrm{rec}}
    -
    \mathcal P_{\mathrm{rec}\rightarrow\mathrm{other}}.
    \label{eq:AT_receiver_loading_ledger}
\end{equation}

Consequently, a small pre-transparency photon-relaxation exposure,
\begin{equation}
    \Gamma_{\mathrm D}t_{\mathrm a}
    \ll
    1,
\end{equation}
does not imply
\begin{equation}
    u_{\mathrm{rec},0}
    \simeq
    0.
    \label{eq:AT_small_exposure_not_cold}
\end{equation}

The initial receiver energy must therefore be obtained independently from the
preceding formation and annealing history.


\section{Conditional Finite Active-to-Latent Interface Specialization}

The general Fossil closure permits a specific TFT specialization associated
with finite-energy active-to-latent dynamics at the FCC--SC transition.

This candidate is distinct from both the ideal Hard Floor and a literal
translation of a single geometric phase boundary.

The relevant distinction is
\begin{equation}
    \begin{aligned}
        \text{bulk FCC--SC conversion coordinate}
        &\neq
        \text{two-dimensional translating phase boundary}
        \\
        &\neq
        \text{finite three-dimensional active-to-latent field}.
    \end{aligned}
    \label{eq:AT_interface_three_way_distinction}
\end{equation}

The first is not generically gapless, while the second possesses
two-dimensional rather than three-dimensional thermodynamic state counting.

The third remains a conditional microscopic candidate.


\subsection{Topology-Normalized Interface Density}

The TFT topology-normalized density law is
\begin{equation}
    \rho_i(r)
    =
    F_i\rho_{\mathrm S}
    \left(
        \frac{\ell_{\mathrm S}}{r}
    \right)^4.
    \label{eq:AT_topology_density_law}
\end{equation}

For the spherical Soft-Floor convergence topology,
\begin{equation}
    F_{\mathrm D}
    =
    \frac{3}{2},
    \label{eq:AT_FD}
\end{equation}
and
\begin{equation}
    \rho_{\mathrm{QG}}
    =
    F_{\mathrm D}
    \rho_{\mathrm S}
    \left(
        \frac{\ell_{\mathrm S}}{r_{\mathrm{QG}}}
    \right)^4.
    \label{eq:AT_rhoQG}
\end{equation}

For a distinct finite interface topology with factor
\(F_{\mathrm I}\),
\begin{equation}
    \rho_{\mathrm I}
    =
    F_{\mathrm I}
    \rho_{\mathrm S}
    \left(
        \frac{\ell_{\mathrm S}}{r_{\mathrm{QG}}}
    \right)^4,
\end{equation}
and therefore
\begin{equation}
    \rho_{\mathrm I}
    =
    \frac{F_{\mathrm I}}{F_{\mathrm D}}
    \rho_{\mathrm{QG}}.
    \label{eq:AT_interface_density_ratio}
\end{equation}

For the conditional rectilinear normal-interface topology,
\begin{equation}
    F_{\mathrm I}
    =
    1,
    \qquad
    \rho_{\mathrm I}
    =
    \frac{\rho_{\mathrm{QG}}}{F_{\mathrm D}}.
    \label{eq:AT_interface_FI1}
\end{equation}

The assignment
\(F_{\mathrm I}=1\)
is a topology hypothesis and is not inferred from the observed CMB
temperature.


\subsection{Conditional Interface Speed}

Assume that the finite active-to-latent field inherits the retained active-FCC
gradient-stiffness scale while carrying the topology-normalized inertia of
Equation~\eqref{eq:AT_interface_density_ratio}.

The corresponding assignments are
\begin{equation}
    K_{\mathrm I}
    =
    \beta_{\mathrm{geo}}
    \rho_{\mathrm D}
    c_{\mathrm E}^2,
    \qquad
    \rho_{\mathrm I}^{(\mathrm{cell})}
    =
    \beta_{\mathrm{geo}}
    \rho_{\mathrm I}.
    \label{eq:AT_interface_stiffness_inertia}
\end{equation}

Here
\(K_{\mathrm I}\)
has dimensions of pressure and
\(\rho_{\mathrm I}^{(\mathrm{cell})}\)
has dimensions of mass density.

The long-wavelength propagation speed is therefore
\begin{equation}
    v_{\mathrm I}^2
    =
    \frac{
        K_{\mathrm I}
    }{
        \rho_{\mathrm I}^{(\mathrm{cell})}
    }
    =
    c_{\mathrm E}^2
    \frac{
        \rho_{\mathrm D}
    }{
        \rho_{\mathrm I}
    }.
    \label{eq:AT_interface_speed_density}
\end{equation}

Using the TFT scale identities,
\begin{equation}
    v_{\mathrm I}(F_{\mathrm I})
    =
    c_{\mathrm E}
    \frac{
        F_{\mathrm D}
    }{
        \sqrt{F_{\mathrm I}}
    }
    \sqrt{
        \frac{
            \ell_{\mathrm S}
        }{
            r_{\mathrm D}
        }
    }.
    \label{eq:AT_interface_speed}
\end{equation}

For
\(F_{\mathrm I}=1\),
\begin{equation}
    v_{\mathrm I}
    =
    c_{\mathrm E}
    F_{\mathrm D}
    \sqrt{
        \frac{
            \ell_{\mathrm S}
        }{
            r_{\mathrm D}
        }
    }
    \simeq
    20.95~{\rm km\,s^{-1}}.
    \label{eq:AT_interface_speed_numeric}
\end{equation}

No measured CMB temperature enters
Equation~\eqref{eq:AT_interface_speed_numeric}.

The unresolved step is the independent microscopic derivation of the
stiffness inheritance in
Equation~\eqref{eq:AT_interface_stiffness_inertia}.


\subsection{Conditional Interface Thermal Capacity}

If the interface candidate is an independent three-dimensional gapless scalar
branch with
\begin{equation}
    \Omega_{\mathrm I}(k_{\mathrm I})
    =
    v_{\mathrm I}k_{\mathrm I},
    \label{eq:AT_interface_gapless_dispersion}
\end{equation}
where \(k_{\mathrm I}\) is the dimensional wave number of the interface
branch, then
\begin{equation}
    u_{\mathrm I}(T)
    =
    a_{\mathrm I}T^4,
    \qquad
    a_{\mathrm I}
    =
    \frac{
        \pi^2k_{\mathrm B}^4
    }{
        30\hbar^3v_{\mathrm I}^3
    }.
    \label{eq:AT_interface_thermal_energy}
\end{equation}

For the electromagnetic field,
\begin{equation}
    a_\gamma
    =
    \frac{
        \pi^2k_{\mathrm B}^4
    }{
        15\hbar^3c_{\mathrm E}^3
    }.
\end{equation}

Hence
\begin{equation}
    \frac{a_{\mathrm I}}{a_\gamma}
    =
    \frac{1}{2}
    \left(
        \frac{
            c_{\mathrm E}
        }{
            v_{\mathrm I}
        }
    \right)^3
    =
    \frac{
        F_{\mathrm I}^{3/2}
    }{
        2F_{\mathrm D}^3
    }
    \left(
        \frac{
            r_{\mathrm D}
        }{
            \ell_{\mathrm S}
        }
    \right)^{3/2}.
    \label{eq:AT_interface_heat_capacity_ratio}
\end{equation}

For
\(F_{\mathrm I}=1\),
the current TFT microscopic quantities give approximately
\begin{equation}
    \frac{a_{\mathrm I}}{a_\gamma}
    \simeq
    1.47\times10^{12}.
    \label{eq:AT_interface_heat_capacity_ratio_numeric}
\end{equation}

This ratio is a conditional thermodynamic state-counting result.

It does not establish the existence, dimensional support, gaplessness,
initial loading, or photon coupling of the proposed branch.


\subsection{Conditional Cold-Interface Endpoint}

Take
\begin{equation}
    T_{\gamma,0}
    =
    T_{\mathrm a},
    \qquad
    T_{\mathrm{rec},0}
    =
    T_{\mathrm I,a}.
\end{equation}

For the limiting condition
\begin{equation}
    T_{\mathrm I,a}
    \rightarrow
    0,
\end{equation}
energy conservation gives
\begin{equation}
    a_\gamma T_{\mathrm a}^4
    =
    \left(
        a_\gamma+a_{\mathrm I}
    \right)
    T_{\mathrm f}^4.
\end{equation}

Therefore
\begin{equation}
    T_{\mathrm f}
    =
    T_{\mathrm a}
    \left[
        1+
        \frac{
            F_{\mathrm I}^{3/2}
        }{
            2F_{\mathrm D}^3
        }
        \left(
            \frac{
                r_{\mathrm D}
            }{
                \ell_{\mathrm S}
            }
        \right)^{3/2}
    \right]^{-1/4}.
    \label{eq:AT_interface_temperature_formula}
\end{equation}

For
\begin{equation}
    F_{\mathrm I}=1,
    \qquad
    F_{\mathrm D}=\frac32,
    \qquad
    T_{\mathrm a}=3000~{\rm K},
\end{equation}
the unrounded TFT microscopic quantities corresponding to
Equation~\eqref{eq:AT_interface_speed_numeric}
give
\begin{equation}
    T_{\mathrm f}
    \simeq
    2.726~{\rm K}.
    \label{eq:AT_interface_temperature_numeric}
\end{equation}

The measured present CMB temperature is not inserted upstream in this
calculation.

Equation~\eqref{eq:AT_interface_temperature_numeric}
is therefore an illustrative conditional consistency result of this
finite-interface specialization.

Its physical realization requires the corresponding interface mode,
three-dimensional state counting, sufficiently small gap, stiffness
normalization, primordial loading, and thermalizing photon interaction.


\subsection{Gap Constraint and the Role of \(r_{\mathrm{QG}}\)}

A general low-wave-number interface branch may instead have
\begin{equation}
    \Omega_{\mathrm I}^2(k_{\mathrm I})
    =
    \Omega_0^2
    +
    v_{\mathrm I}^2k_{\mathrm I}^2.
    \label{eq:AT_interface_gapped_dispersion}
\end{equation}

Its scalar thermal energy density can be written
\begin{equation}
    u_{\mathrm I}(T)
    =
    a_{\mathrm I}T^4
    \mathcal G
    \left(
        \frac{
            \hbar\Omega_0
        }{
            k_{\mathrm B}T
        }
    \right),
    \label{eq:AT_interface_gapped_energy}
\end{equation}
where
\begin{equation}
    \mathcal G(x)
    =
    \frac{15}{\pi^4}
    \int_0^\infty
    \frac{
        y^2\sqrt{y^2+x^2}
    }{
        \exp
        \left(
            \sqrt{y^2+x^2}
        \right)-1
    }
    \,\mathrm dy,
    \qquad
    \mathcal G(0)=1.
    \label{eq:AT_interface_gap_function}
\end{equation}

The gapless approximation requires
\begin{equation}
    \Omega_0
    \ll
    \frac{k_{\mathrm B}T}{\hbar}
    \label{eq:AT_interface_gap_requirement}
\end{equation}
over the thermally relevant regime.

At a few kelvin,
\begin{equation}
    \frac{k_{\mathrm B}T}{\hbar}
    \sim
    10^{11}~{\rm s^{-1}}.
\end{equation}

The microscopic convergence radius
\begin{equation}
    r_{\mathrm{QG}}
    \sim
    10^{-29}~{\rm m}
\end{equation}
can enter the interface density and inertia normalization, but it cannot by
itself provide the required infrared frequency scale because
\begin{equation}
    \frac{c_{\mathrm E}}{r_{\mathrm{QG}}}
    \sim
    10^{36}~{\rm s^{-1}}.
    \label{eq:AT_rQG_frequency_scale}
\end{equation}

Thus the role of
\(r_{\mathrm{QG}}\)
in the present construction is mechanical normalization rather than direct
generation of the low-frequency thermal scale.

A sufficiently small
\(\Omega_0\)
must arise from an independently derived collective soft mode, symmetry,
critical structure, or another genuine infrared mechanism.


\subsection{Energy-Conserving Two-Sector Evolution}

Let
\begin{equation}
    \mathcal P_{\gamma\mathrm{rec}}
    \left(
        T_\gamma,
        T_{\mathrm{rec}}
    \right)
\end{equation}
denote the net microscopic power transferred per unit physical volume from the
thermal photon sector into the receiver, with positive sign denoting energy
flow out of the photons.

The coupled temperature equations are
\begin{equation}
    C_V^{(\gamma)}
    \left(
        T_\gamma
    \right)
    \dot T_\gamma
    =
    -
    \mathcal P_{\gamma\mathrm{rec}}
    \left(
        T_\gamma,
        T_{\mathrm{rec}}
    \right),
    \label{eq:AT_photon_temperature_evolution}
\end{equation}
and
\begin{equation}
    C_V^{(\mathrm{rec})}
    \left(
        T_{\mathrm{rec}}
    \right)
    \dot T_{\mathrm{rec}}
    =
    +
    \mathcal P_{\gamma\mathrm{rec}}
    \left(
        T_\gamma,
        T_{\mathrm{rec}}
    \right).
    \label{eq:AT_receiver_temperature_evolution}
\end{equation}

For the electromagnetic field,
\begin{equation}
    C_V^{(\gamma)}(T)
    =
    \frac{
        4\pi^2k_{\mathrm B}^4
    }{
        15\hbar^3c_{\mathrm E}^3
    }
    T^3,
    \label{eq:AT_photon_heat_capacity}
\end{equation}
while for the selected linear receiving branches,
\begin{equation}
    C_V^{(\mathrm{rec})}(T)
    =
    \frac{
        2\pi^2k_{\mathrm B}^4
    }{
        15\hbar^3
    }
    \left(
        \sum_jv_j^{-3}
    \right)
    T^3.
    \label{eq:AT_receiver_heat_capacity}
\end{equation}

Adding the two equations gives
\begin{equation}
    \frac{\mathrm d}{\mathrm dt}
    \left[
        u_{\mathrm{CMB}}(T_\gamma)
        +
        u_{\mathrm{rec}}(T_{\mathrm{rec}})
    \right]
    =
    0.
    \label{eq:AT_two_sector_energy_conservation}
\end{equation}

The microscopic transfer kernel therefore determines the rate of approach to
equilibrium, while the endpoint is fixed by the conserved total energy once
the initial states and receiver thermodynamics have been specified.


\subsection{Logical Status of the Fossil Closure}

The general Fossil result has the conditional structure
\begin{equation}
    \left\{
        T_{\gamma,0},
        \,
        T_{\mathrm{rec},0},
        \,
        \text{receiving spectrum}
    \right\}
    \longrightarrow
    T_{\mathrm{eq}}.
    \label{eq:AT_fossil_conditional_status}
\end{equation}

The thermodynamic density of states does not independently determine
\begin{equation}
    T_{\mathrm{rec},0},
    \qquad
    \mathcal P_{\gamma\mathrm{rec}},
\end{equation}
and the carrier-projected dissipative spectrum does not determine the
thermodynamic multiplicity.

The strongest general result of the closed Fossil construction is the
energy-conserving two-sector endpoint in
Equations~\eqref{eq:AT_fossil_energy_conservation_general}
and~\eqref{eq:AT_fossil_temperature_general}.

The finite active-to-latent interface provides the illustrative conditional
specialization
\begin{equation}
    \begin{aligned}
        F_{\mathrm I}=1,
        \quad
        &\text{independent three-dimensional gapless scalar receiver},
        \\
        T_{\mathrm I,a}\simeq0
        \quad
        &\Longrightarrow
        \quad
        T_{\mathrm f}\simeq2.726~{\rm K}.
    \end{aligned}
    \label{eq:AT_interface_conditional_status}
\end{equation}

The observed CMB temperature is not used to infer
\(F_{\mathrm I}\),
\(v_{\mathrm I}\),
\(\Omega_0\),
or
\(T_{\mathrm I,a}\).

Accordingly, the general Fossil endpoint is the retained energy-conservation
result for a specified microscopic completion, while the finite-interface
calculation serves as a downstream consistency test of one particular TFT
realization.


\section{Macroscopic Thermal Energy Balance}

The Fossil-radiation closure derived above describes a closed local
photon--receiver subsystem during an interval in which external work,
macroscopic volume change, and exchange with additional sectors are
negligible.

The macroscopic thermal environment is more general.

Once continuous energy injection, terminal conversion, transport into
additional Active-Vacuum degrees of freedom, or capture by State-II
structures is admitted, the thermalized sector must be treated as an open
subsystem.

Let
\(u_{\mathrm{th}}(T)\)
denote the total energy density assigned to the specified thermalized sector,
and define its volumetric heat capacity by
\begin{equation}
    C_V^{(\mathrm{th})}(T)
    \equiv
    \frac{
        \mathrm d u_{\mathrm{th}}
    }{
        \mathrm dT
    }.
    \label{eq:AT_heat_capacity_definition}
\end{equation}

For a spatially homogeneous coarse-graining volume in which mechanical work
and explicit spatial transport may be neglected over the interval of
interest, the local thermal balance is
\begin{equation}
    C_V^{(\mathrm{th})}(T)
    \frac{
        \mathrm dT
    }{
        \mathrm dt
    }
    =
    \mathcal P_{\mathrm{in}}(T)
    -
    \mathcal P_{\mathrm{sink}}(T).
    \label{eq:AT_thermal_balance}
\end{equation}

The quantities
\(\mathcal P_{\mathrm{in}}\)
and
\(\mathcal P_{\mathrm{sink}}\)
are power densities with dimensions
\(\mathrm{W\,m^{-3}}\).

If spatial heat transport, macroscopic compression or expansion, or
mechanical work becomes important, the corresponding flux-divergence and
work terms must be restored before
Equation~\eqref{eq:AT_thermal_balance}
is applied.

The thermal input may be decomposed as
\begin{equation}
    \mathcal P_{\mathrm{in}}
    =
    \mathcal P_{\mathrm{cont}}
    +
    \mathcal P_{\mathrm{term}}
    +
    \mathcal P_{\mathrm{other}}.
    \label{eq:AT_input_power}
\end{equation}

Here
\(\mathcal P_{\mathrm{cont}}\)
denotes thermal power originating from energy transferred continuously out of
the coherent photon population,
\(\mathcal P_{\mathrm{term}}\)
denotes thermal power originating from terminal coherence-loss events, and
\(\mathcal P_{\mathrm{other}}\)
collects any additional independently derived thermal source.

The distinction between energy leaving the coherent carrier and energy
entering the thermalized sector is essential.

Let
\(u_\gamma^{(\mathrm{coh})}\)
denote the local energy density carried by the coherent photon population to
which the inherited scale-transfer law applies.

The instantaneous rate at which energy leaves that coherent population is
\begin{equation}
    \mathcal P_{\mathrm{loss}}^{(\mathrm{coh})}
    =
    \Gamma_{\mathrm D}
    u_\gamma^{(\mathrm{coh})}.
    \label{eq:AT_coherent_loss_power_density}
\end{equation}

Equation~\eqref{eq:AT_coherent_loss_power_density}
is a statement about the coherent carrier sector.

It does not imply that the same power enters the thermalized sector at the
same instant.

Transferred coherent energy may first populate nonthermal receiving modes,
intermediate collective excitations, or another explicitly tracked
finite-energy sector and thermalize only after a finite delay.

Consequently, no general pointwise equality is imposed between
\(\mathcal P_{\mathrm{cont}}\)
and
\(\mathcal P_{\mathrm{loss}}^{(\mathrm{coh})}\).

In the special local and effectively memoryless limit in which continuously
transferred coherent energy thermalizes without appreciable intermediate
storage on the temporal and spatial coarse-graining scales being modeled,
one obtains
\begin{equation}
    \mathcal P_{\mathrm{cont}}
    =
    \Gamma_{\mathrm D}
    u_\gamma^{(\mathrm{coh})}.
    \label{eq:AT_complete_continuous_thermalization}
\end{equation}

Outside this limit,
\(\mathcal P_{\mathrm{cont}}(t)\)
depends on the preceding population and relaxation of the intermediate
receiving sector and need not coincide instantaneously with the current
coherent loss power.

The general causal structure is therefore
\begin{equation}
    \text{coherent energy loss}
    \longrightarrow
    \text{nonthermal receiving sector}
    \longrightarrow
    \text{thermal input}.
    \label{eq:AT_coherent_to_thermal_sequence}
\end{equation}

The intermediate stage may collapse to an effectively instantaneous process
only when its relaxation timescale is negligible relative to the retained
coarse-graining timescale.

Terminal coherence-loss events require the same distinction.

Let
\(\mathrm d\dot n_{\mathrm{term}}\)
denote the differential terminal-event rate per physical volume.

The gross power released from the coherent terminal population is
\begin{equation}
    \mathcal P_{\mathrm{release}}^{(\mathrm{term})}
    =
    \int
    E_{\mathrm{term}}
    \,\mathrm d\dot n_{\mathrm{term}}.
    \label{eq:AT_terminal_release_power_density}
\end{equation}

If this released energy is transferred directly into the thermalized sector
on the retained coarse-graining timescale, then
\begin{equation}
    \mathcal P_{\mathrm{term}}
    =
    \mathcal P_{\mathrm{release}}^{(\mathrm{term})}
    =
    \int
    E_{\mathrm{term}}
    \,\mathrm d\dot n_{\mathrm{term}}.
    \label{eq:AT_terminal_power_density}
\end{equation}

If instead the terminal energy first occupies a nonthermal intermediate
sector, its eventual thermal contribution must be propagated through that
sector rather than identified instantaneously with
Equation~\eqref{eq:AT_terminal_release_power_density}.

Thus both continuous and terminal coherent-energy loss obey the same
thermodynamic principle:
\begin{equation}
    \text{energy leaving a coherent carrier}
    \not\equiv
    \text{instantaneous thermal heating}.
    \label{eq:AT_loss_not_instantaneous_heating}
\end{equation}

A stationary thermal state,
denoted by
\(T_{\mathrm{fp}}\),
satisfies
\begin{equation}
    \mathcal P_{\mathrm{in}}
    \left(
        T_{\mathrm{fp}}
    \right)
    =
    \mathcal P_{\mathrm{sink}}
    \left(
        T_{\mathrm{fp}}
    \right).
    \label{eq:AT_fixed_point}
\end{equation}

The subscript ``fp'' denotes a driven thermal fixed point and distinguishes
this quantity from the closed-system Fossil equilibrium temperature
\(T_{\mathrm{eq}}\).

Linear stability requires
\begin{equation}
    \left.
    \frac{
        \mathrm d
    }{
        \mathrm dT
    }
    \left[
        \mathcal P_{\mathrm{in}}(T)
        -
        \mathcal P_{\mathrm{sink}}(T)
    \right]
    \right|_{T_{\mathrm{fp}}}
    <
    0.
    \label{eq:AT_fixed_point_stability}
\end{equation}

Equivalently,
\begin{equation}
    \left.
    \frac{
        \mathrm d\mathcal P_{\mathrm{sink}}
    }{
        \mathrm dT
    }
    \right|_{T_{\mathrm{fp}}}
    >
    \left.
    \frac{
        \mathrm d\mathcal P_{\mathrm{in}}
    }{
        \mathrm dT
    }
    \right|_{T_{\mathrm{fp}}}.
    \label{eq:AT_sink_slope_condition}
\end{equation}

The two thermal constructions therefore represent different physical limits:
\begin{equation}
    \begin{aligned}
        T_{\mathrm{eq}}:
        &\quad
        \text{closed photon--receiver energy redistribution},
        \\
        T_{\mathrm{fp}}:
        &\quad
        \text{open driven source--sink stationary state}.
    \end{aligned}
    \label{eq:AT_equilibrium_fixed_point_distinction}
\end{equation}

No equality between
\(T_{\mathrm{eq}}\)
and
\(T_{\mathrm{fp}}\)
is assumed.


\section{Candidate State-II Thermal Sink}

Localized FCC-to-State-II phase structures provide a structurally motivated
candidate sink for thermal Active-Vacuum energy.

This open-system sink is distinct from the finite active-to-latent receiving
sector used in the closed Fossil-equilibrium construction.

The present construction does not assume that State-II structures are perfect
absorbers.

Let
\(\sigma_i^{(\mathrm{abs})}(\Omega)\)
denote the effective absorption cross section of State-II structure \(i\) for
a thermal excitation of angular frequency
\(\Omega\).

Define the corresponding inverse effective absorption length by
\begin{equation}
    \ell_{\mathrm{II}}^{-1}(\Omega)
    \equiv
    \frac{1}{V}
    \sum_i
    \sigma_i^{(\mathrm{abs})}(\Omega).
    \label{eq:AT_StateII_inverse_absorption_length}
\end{equation}

Its dimensions are
\begin{equation}
    \left[
        \ell_{\mathrm{II}}^{-1}
    \right]
    =
    \mathrm{m^{-1}}.
\end{equation}

Let
\(u_\Omega(T)\)
denote the thermal spectral energy density per unit angular frequency of the
excitations participating in the capture process, and let
\(v(\Omega)\)
denote their propagation speed.

The State-II capture power density is then
\begin{equation}
    \mathcal P_{\mathrm{II}}(T)
    =
    \int_0^\infty
    v(\Omega)
    \ell_{\mathrm{II}}^{-1}(\Omega)
    u_\Omega(T)
    \,\mathrm d\Omega.
    \label{eq:AT_StateII_sink}
\end{equation}

Equation~\eqref{eq:AT_StateII_sink}
has dimensions
\(\mathrm{W\,m^{-3}}\).

The sink strength therefore depends on the State-II population, effective
capture geometry, and frequency-dependent interaction with the participating
thermal modes rather than on State-II mass alone.

The function
\(\ell_{\mathrm{II}}^{-1}(\Omega)\)
is not presently derived from the microscopic FCC-to-State-II phase-boundary
dynamics.

Accordingly, State-II structures constitute a candidate sink architecture,
not a completed thermal sink law.

In a gray-capture approximation over the thermally relevant frequency range,
take
\begin{equation}
    \ell_{\mathrm{II}}^{-1}(\Omega)
    \simeq
    \ell_{\mathrm{II}}^{-1},
    \qquad
    v(\Omega)
    \simeq
    v_{\mathrm{th}}.
    \label{eq:AT_gray_capture_approximation}
\end{equation}

Define the corresponding effective capture rate
\begin{equation}
    \Gamma_{\mathrm{II}}^{(\mathrm{cap})}
    \equiv
    v_{\mathrm{th}}
    \ell_{\mathrm{II}}^{-1}.
    \label{eq:AT_StateII_capture_rate}
\end{equation}

Then
\begin{equation}
    \mathcal P_{\mathrm{II}}
    =
    \Gamma_{\mathrm{II}}^{(\mathrm{cap})}
    u_{\mathrm{th}}.
    \label{eq:AT_gray_sink}
\end{equation}

The gray approximation does not derive either
\(v_{\mathrm{th}}\)
or
\(\ell_{\mathrm{II}}\);
it only reduces the frequency-dependent capture problem to an effective rate.

The SC State-II configuration is not interpreted as an ordinary thermal
reservoir.

Energy captured from the thermalized Active Vacuum must ultimately appear as
an explicitly identified change in structural energy, phase-boundary work,
defect growth, or another finite-energy mechanical degree of freedom.

No ordinary finite thermal inertia is assigned to the zero-inertia Hard
Floor.

Thus,
\begin{equation}
    \text{State-II capture}
    \neq
    \text{Hard-Floor thermalization}.
    \label{eq:AT_StateII_HardFloor_distinction}
\end{equation}

\section{Absolute-Temperature Normalization Test}

A thermal fixed-point equation becomes an absolute temperature prediction only
if every quantity controlling its source, sink, and thermodynamic
normalization is determined independently of the observed temperature.

Suppose, for illustration, that the specified thermalized sector is described
over the relevant temperature range by independent three-dimensional
gapless linear bosonic branches indexed by \(j\).

Their total thermal energy density is
\begin{equation}
    u_{\mathrm{th}}(T)
    =
    \frac{
        \pi^2k_{\mathrm B}^4
    }{
        30\hbar^3
    }
    \left(
        \sum_j
        v_j^{-3}
    \right)
    T^4,
    \label{eq:AT_total_linear_thermal_energy}
\end{equation}
where each independent physical polarization or branch is counted separately.

If the electromagnetic field is included in the sum, its two transverse
polarizations each contribute a branch with speed
\(c_{\mathrm E}\).

The corresponding volumetric heat capacity is
\begin{equation}
    C_V^{(\mathrm{th})}(T)
    =
    \frac{
        2\pi^2k_{\mathrm B}^4
    }{
        15\hbar^3
    }
    \left(
        \sum_j
        v_j^{-3}
    \right)
    T^3.
    \label{eq:AT_total_linear_heat_capacity}
\end{equation}

These expressions apply only within the low-temperature linear portions of
the specified branches. If finite-band or nonlinear-dispersion effects are
thermally sampled, the exact density of states must replace the linear
continuum approximation.

Now consider the gray State-II sink of
Equation~\eqref{eq:AT_gray_sink}.

Over an interval in which the thermal input
\begin{equation}
    \mathcal P_0
    \equiv
    \mathcal P_{\mathrm{cont}}
    +
    \mathcal P_{\mathrm{term}}
    +
    \mathcal P_{\mathrm{other}}
    \label{eq:AT_constant_thermal_input}
\end{equation}
may be treated as approximately independent of \(T\),
Equation~\eqref{eq:AT_thermal_balance}
becomes
\begin{equation}
    \frac{
        2\pi^2k_{\mathrm B}^4
    }{
        15\hbar^3
    }
    \left(
        \sum_j
        v_j^{-3}
    \right)
    T^3\dot T
    =
    \mathcal P_0
    -
    \Gamma_{\mathrm{II}}^{(\mathrm{cap})}
    \frac{
        \pi^2k_{\mathrm B}^4
    }{
        30\hbar^3
    }
    \left(
        \sum_j
        v_j^{-3}
    \right)
    T^4.
    \label{eq:AT_gray_temperature_evolution}
\end{equation}

The positive stationary solution is
\begin{equation}
    T_{\mathrm{fp}}
    =
    \left[
        \frac{
            \mathcal P_0
        }{
            \Gamma_{\mathrm{II}}^{(\mathrm{cap})}
            \dfrac{
                \pi^2k_{\mathrm B}^4
            }{
                30\hbar^3
            }
            \left(
                \displaystyle\sum_j
                v_j^{-3}
            \right)
        }
    \right]^{1/4}.
    \label{eq:AT_temperature_fixed_point}
\end{equation}

If complete thermalization of continuously transferred coherent energy is
independently established over the retained coarse-graining interval, one may
specialize
\begin{equation}
    \mathcal P_{\mathrm{cont}}
    =
    \Gamma_{\mathrm D}
    u_\gamma^{(\mathrm{coh})}.
    \label{eq:AT_complete_thermalization_specialization}
\end{equation}

Without that additional condition,
Equation~\eqref{eq:AT_complete_thermalization_specialization}
is not imposed automatically.

For constant
\(\mathcal P_0\)
and positive
\(\Gamma_{\mathrm{II}}^{(\mathrm{cap})}\),
linearization around
\(T_{\mathrm{fp}}\)
gives
\begin{equation}
    \delta\dot T
    =
    -
    \Gamma_{\mathrm{II}}^{(\mathrm{cap})}
    \delta T.
    \label{eq:AT_temperature_stability}
\end{equation}

Thus the gray fixed point is linearly stable within the assumptions of the
reduced source--sink model.

Its numerical evaluation requires independent determination of
\begin{equation}
    \mathcal P_{\mathrm{cont}},
    \qquad
    \mathcal P_{\mathrm{term}},
    \qquad
    \mathcal P_{\mathrm{other}},
    \qquad
    \Gamma_{\mathrm{II}}^{(\mathrm{cap})},
    \qquad
    \{v_j\}.
    \label{eq:AT_open_temperature_parameters}
\end{equation}

The branch speeds
\(\{v_j\}\)
determine the thermal state-counting normalization, while
\(\Gamma_{\mathrm{II}}^{(\mathrm{cap})}\)
sets the gray sink normalization.

The source normalization is independent of both.

The observed CMB temperature therefore remains a downstream comparison rather
than an upstream condition used to determine any quantity in
Equation~\eqref{eq:AT_open_temperature_parameters}.

In particular,
\begin{equation}
    R_{\mathrm U}
    =
    \frac{
        c_{\mathrm E}
    }{
        \Gamma_{\mathrm D}
    }
\end{equation}
does not by itself supply the thermal energy-density normalization.

More generally, dimensional combinations of
\(\Gamma_{\mathrm D}\),
\(R_{\mathrm U}\),
\(c_{\mathrm E}\),
\(\hbar\),
and
\(k_{\mathrm B}\)
constitute a physical temperature prediction only when a microscopic
mechanism selects the corresponding combination.

A second normalization result follows if both input and sink powers are
proportional to the same thermal energy density.

Suppose
\begin{equation}
    \mathcal P_{\mathrm{in}}
    =
    \Gamma_{\mathrm{in}}
    u_{\mathrm{th}},
    \qquad
    \mathcal P_{\mathrm{sink}}
    =
    \Gamma_{\mathrm{sink}}
    u_{\mathrm{th}}.
    \label{eq:AT_proportional_rates}
\end{equation}

For a nonzero thermal state, stationarity requires only
\begin{equation}
    \Gamma_{\mathrm{in}}
    =
    \Gamma_{\mathrm{sink}}.
    \label{eq:AT_proportional_rate_equality}
\end{equation}

The temperature cancels.

Therefore,
\begin{equation}
    \mathcal P_{\mathrm{in}}
    \propto
    u_{\mathrm{th}},
    \qquad
    \mathcal P_{\mathrm{sink}}
    \propto
    u_{\mathrm{th}}
    \not\Rightarrow
    \text{a unique nonzero thermal temperature}.
    \label{eq:AT_temperature_normalization_no_go}
\end{equation}

A unique absolute temperature requires an independently normalized energy
scale, state-counting coefficient, nonlinear source or sink dependence, or
boundary condition that survives the stationary balance.

The same principle applies to the closed Fossil relation in
Equation~\eqref{eq:AT_fossil_temperature_general}:
energy conservation determines
\(T_{\mathrm{eq}}\)
once the initial photon state, initial receiver state, and receiver
thermodynamics have been independently specified.

The conditional interface result in
Equation~\eqref{eq:AT_interface_temperature_numeric}
is a particular realization of this general structure.

Its numerical value becomes a physical prediction if the interface topology,
three-dimensional mode support, stiffness, gap, primordial loading, and
thermalizing interaction are independently established.

Accordingly, the general closed Fossil equilibrium equation, the
finite-interface specialization, and the open driven fixed-point equation
provide complementary routes toward an absolute CMB temperature, with the
remaining normalization supplied by their corresponding microscopic
completion.


\section{Temperature Evolution Seen by Local Matter}

The local thermodynamic state of the radiation field and the temperature
inferred by an astronomical matter system are related but distinct
quantities.

The thermal equations above describe the state of a specified local radiation
and Active-Vacuum subsystem.

A temperature measured at another event additionally depends on propagation
history, relative motion, endpoint gravitational loading, angular radiation
distribution, and the matter--radiation interaction through which the
spectrum is sampled.

For stationary gravitational conditions and the factorized propagation
regime, Equation~\eqref{eq:AT_redshift_hierarchy} gives
\begin{equation}
    \boxed{
    1+z_{\mathrm{obs}}
    =
    (1+z_{\mathrm{kin}})
    (1+z_{\mathrm{grav}})
    (1+z_{\mathrm D}).
    }
\end{equation}

The gravity-neutral endpoint limit is recovered when
\begin{equation}
    z_{\mathrm{grav}}
    \rightarrow
    0.
\end{equation}

A completed TFT CMB construction must therefore derive a local
matter--radiation measurement relation of the schematic form
\begin{equation}
    T_{\mathrm{meas}}
    =
    T_{\mathrm{rad}}
    \left(
        \mathbf{x},
        t,
        v_{\mathrm{ext}},
        \alpha_{\mathrm G},
        \text{local radiation state},
        \text{local matter response}
    \right).
    \label{eq:AT_temperature_observer_map}
\end{equation}

Equation~\eqref{eq:AT_temperature_observer_map}
specifies the required dependencies; it is not itself a constitutive
temperature law.

Relative motion through an isotropic thermal field generally produces an
angularly dependent observed distribution rather than a single
direction-independent temperature obtained by multiplying the rest-frame
temperature by a one-dimensional Doppler factor.

Therefore,
\begin{equation}
    T_{\mathrm{meas}}
    =
    T_{\mathrm{meas}}
    \left(
        z_{\mathrm{obs}}
    \right)
\end{equation}
cannot be assumed as an upstream thermal closure.

The complete observer construction must distinguish the local thermodynamic
radiation state, the propagation contribution
\(z_{\mathrm D}\),
the relative-motion contribution
\(z_{\mathrm{kin}}\),
the endpoint gravitational contribution
\(z_{\mathrm{grav}}\),
the angular radiation distribution in the matter frame, and the microscopic
matter--radiation response.

Consequently,
\begin{equation}
    T_{\mathrm{meas}}(z_{\mathrm{obs}})
    \text{ is a downstream observational map,
    not an upstream thermal closure.}
    \label{eq:AT_temperature_redshift_status}
\end{equation}

The present appendix does not derive the complete observer map.

Measurements of CMB temperature at nonzero astronomical redshift therefore
remain independent empirical tests of the completed TFT thermal and
propagation sectors.


\section{Structural Constraints and Logical Status}

The preceding construction defines a thermodynamically consistent framework
for testing microscopic CMB completions.

The following constraints distinguish general results established by the
thermal analysis from properties specific to candidate receiving mechanisms.


\subsection{Coherent Propagation and Thermal Equilibration}

The coherent scale-transfer branch preserves the identity and photon number
of the surviving carrier, whereas a zero-chemical-potential Planck state
requires a distinct photon-number-changing thermal interaction.

Therefore,
\begin{equation}
    \text{coherent scale transfer}
    \not\Rightarrow
    \text{thermal Planck evolution}.
    \label{eq:AT_constraint_scale_only}
\end{equation}

The thermal completion requires microscopic detailed balance together with
an interaction satisfying
\begin{equation}
    \left[
        \hat H_{\mathrm{th}},
        \hat N_{\mathrm{ph}}
    \right]
    \neq
    0.
\end{equation}


\subsection{Finite Modes, Bulk Conversion, and Interface Dimensionality}

A finite set of conservative undamped modes can exchange energy reversibly
with a coherent photon but does not by itself supply the irreversible
continuum required for a nonzero constant low-frequency relaxation rate.

Thus,
\begin{equation}
    \text{finite conservative mode set}
    \not\Rightarrow
    \text{irreversible receiving continuum}.
    \label{eq:AT_constraint_finite_receiver}
\end{equation}

A bulk coordinate distinguishing inequivalent stable structural phases is
also not generically gapless.

Near a stable configuration,
\begin{equation}
    U_{\mathrm{conv}}(\xi)
    =
    U_{\mathrm{conv}}(\xi_0)
    +
    \frac{1}{2}
    K_{\mathrm{conv}}
    \left(
        \xi-\xi_0
    \right)^2
    +
    \cdots,
    \qquad
    K_{\mathrm{conv}}>0.
    \label{eq:AT_conversion_potential}
\end{equation}

For finite generalized inertia
\(I_\xi>0\),
\begin{equation}
    \Omega_{\mathrm{conv},0}^{\,2}
    =
    \frac{
        K_{\mathrm{conv}}
    }{
        I_\xi
    }
    >
    0.
    \label{eq:AT_conversion_gap}
\end{equation}

A bulk FCC--SC conversion coordinate therefore does not automatically provide
the gapless branch required by the massless receiver approximation.

A literal translating FCC--SC phase boundary is different. Its translation
may be gapless, but it is a two-dimensional interface excitation.

For a two-dimensional gapless linear branch,
\begin{equation}
    \frac{
        \mathcal D_{\mathrm{wall}}(\Omega)
    }{
        \mathcal A_{\mathrm{wall}}
    }
    \propto
    \Omega,
\end{equation}
and its thermal energy per physical interface area scales as
\begin{equation}
    \frac{
        U_{\mathrm{wall}}
    }{
        \mathcal A_{\mathrm{wall}}
    }
    \propto
    T^3.
    \label{eq:AT_wall_thermal_scaling}
\end{equation}

A volumetric contribution then requires an independently derived
interface-area density
\(
\mathcal A_{\mathrm{wall}}/V
\).

The three-dimensional interface specialization developed above is therefore
not identified with a single geometrical wall translation. Its volumetric
state-counting interpretation refers instead to the proposed finite
three-dimensional active-to-latent field.


\subsection{Thermodynamic State Counting versus Irreversible Coupling}

The central microscopic separation established in this appendix is
\begin{equation}
    \mathcal D_{\mathrm{AV}}(\Omega)
    \not\equiv
    \mathcal S_{\gamma\mathrm{AV}}(\Omega).
    \label{eq:AT_constraint_DOS_projection}
\end{equation}

A three-dimensional gapless linear branch may possess
\begin{equation}
    u(T)\propto T^4
\end{equation}
without supplying the Ohmic carrier projection required for a constant
low-frequency damping coefficient.

The reduced FCC internal phase sector illustrates this distinction.

Because its retained contact mechanics depend on phase differences, the
reduced phase action admits a uniform phase-shift symmetry. If a finite
kinetic metric exists, this provides a conditional route to a
three-dimensional gapless collective mode.

However, for the regular local derivative photon--phase couplings examined
in the microscopic analysis, canonical normalization gives the infrared
scaling
\begin{equation}
    \left|
        \langle f|
        H_{\gamma\phi}
        |i\rangle
    \right|^2
    \propto
    \Omega,
\end{equation}
while
\begin{equation}
    \frac{
        \mathcal D_\phi(\Omega)
    }{
        V
    }
    \propto
    \Omega^2.
\end{equation}

Hence
\begin{equation}
    \mathcal S_{\gamma\phi}(\Omega)
    \propto
    \Omega^3
\end{equation}
for this regular local derivative class, rather than the required Ohmic
scaling proportional to
\(\Omega\).

Therefore,
\begin{equation}
    T^4\text{ thermal energy density}
    \not\Rightarrow
    \text{Ohmic photon projection}.
    \label{eq:AT_constraint_T4_Ohmic}
\end{equation}

This leaves open independently derived collective, critical, nonlocal, or
other nonanalytic couplings.


\subsection{Extended Continua and Infrared Scales}

A semi-infinite finite-inertia continuum with outgoing-wave boundary
conditions can conditionally produce an Ohmic endpoint response.

This dynamical result does not establish the required thermodynamic state
counting.

A literal one-dimensional linear bosonic continuum has
\begin{equation}
    \frac{U_{\mathrm{1D}}}{L}
    \propto
    T^2,
\end{equation}
rather than a three-dimensional
\(T^4\)
energy-density law.

Moreover, the scale-depth coordinate appearing in the endpoint construction
is not postulated as an additional ordinary Euclidean spatial dimension.

Therefore,
\begin{equation}
    \text{Ohmic endpoint dynamics}
    \not\Rightarrow
    \text{a derived three-dimensional thermal reservoir}.
    \label{eq:AT_constraint_extended_continuum}
\end{equation}

A separate constraint applies to the origin of the required infrared scale.

Regular finite-dimensional harmonic hybridization can produce mode mixing
and avoided crossings but does not by itself generate a parametrically new
low-frequency scale.

Thus,
\begin{equation}
    \text{regular finite-dimensional harmonic hybridization}
    \not\Rightarrow
    \text{an independently generated infrared scale}.
    \label{eq:AT_constraint_hybridization}
\end{equation}

The interface result derived above is consistent with this restriction:
\(r_{\mathrm{QG}}\)
may normalize the candidate interface density and inertia, but
\(c_{\mathrm E}/r_{\mathrm{QG}}\)
is far too large to supply the required thermal infrared scale.

A sufficiently soft receiver must therefore obtain its low-frequency
structure from a collective mechanism, symmetry, criticality, or another
independently established infrared origin.


\subsection{Energy Accounting and Primordial Loading}

The coherent energy ledger requires continuously transferred energy and any
terminal remainder to remain in the complete accounting.

When
\begin{equation}
    E_{\mathrm{emit}}
    \gg
    E_{\mathrm{term}},
\end{equation}
the continuously transferred contribution dominates.

Therefore,
\begin{equation}
    \text{terminal coherence loss alone}
    \not\Rightarrow
    \text{an energy-complete thermal source}.
    \label{eq:AT_constraint_terminal_only}
\end{equation}

The Fossil equilibrium relation likewise requires an independently specified
initial receiving state.

As established in
Equation~\eqref{eq:AT_small_exposure_not_cold},
a small pre-transparency relaxation exposure,
\begin{equation}
    \Gamma_{\mathrm D}t_{\mathrm a}
    \ll
    1,
\end{equation}
does not imply
\begin{equation}
    u_{\mathrm{rec},0}
    \simeq
    0.
\end{equation}

The actual receiver loading must follow from the catastrophic-decompression,
formation, and annealing history.


\subsection{Absolute-Temperature Constraint}

Neither energy conservation nor the existence of a stable thermal state alone
sets the absolute temperature without specifying the physical state being
equilibrated.

For the closed Fossil system,
\(T_{\mathrm{eq}}\)
is determined once the initial photon state, initial receiver state, and
receiver thermodynamics are specified.

For the finite-interface specialization, the numerical endpoint additionally
depends on the topology, dimensional support, stiffness, gap, primordial
loading, and photon coupling.

For the open driven system,
\(T_{\mathrm{fp}}\)
is determined once the source, thermodynamic normalization, and sink law are
specified.

Consequently,
\begin{equation}
    \text{energy conservation}
    +
    \text{thermal stability}
    \not\Rightarrow
    \text{a unique absolute temperature without microscopic normalization}.
    \label{eq:AT_constraint_absolute_temperature}
\end{equation}

The observed CMB temperature remains the downstream empirical comparison.


\section{Fossil-CMB Predictions, Completion, and Falsification}
\label{sec:AV_predictions_falsification}

The cosmic microwave background presents one of the most
demanding tests of TFT's cosmic architecture. In the Fossil
Route, the observed radiation is inherited from an earlier
epoch of catastrophic decompression and thermal evolution.
Its present spectrum and temperature must therefore emerge
from the physical history of the radiation, not from a
continuously imposed cosmic equilibrium.

A successful mechanical explanation must connect the
primordial radiation, its interaction with the Active Vacuum,
and its subsequent propagation through one consistent
thermodynamic history.

\subsection*{TFT Predictions: The Thermal Endpoint}

For an isolated photon--receiver system of fixed volume,
with independent three-dimensional gapless receiving modes
of linear dispersion, conservation of energy establishes
the conditional equilibrium temperature
%
\begin{equation}
T_{\mathrm{eq}}^4
=
\frac{
2c_{\mathrm E}^{-3}T_{\gamma,0}^4
+
\left(\sum_j v_j^{-3}\right)T_{\mathrm{rec},0}^4
}{
2c_{\mathrm E}^{-3}
+
\sum_j v_j^{-3}
}.
\end{equation}
%
This relation, developed in
Equation~\eqref{eq:AT_fossil_temperature_general},
applies when the thermally populated receiving modes
lie within their assumed linear-dispersion regime.

The positive energy weights impose a definite physical
restriction: the final equilibrium temperature must
lie between the initial photon and receiver temperatures.
If both begin at the same temperature, their equilibration
cannot produce net cooling.

The result expresses an energy-conserving thermal endpoint,
not an independent prediction of today's CMB temperature.
The receiving spectrum, its initial loading, and the
conditions permitting equilibrium must be established
through the microscopic constitution of the Active Vacuum.

A further distinction concerns the blackbody spectrum.
If an initially nonequilibrium photon population has an
incompatible photon number, ordinary number-conserving
scattering cannot generically establish a Planck
distribution with zero chemical potential. The required
thermal mechanism must supply an appropriate
number-changing interaction or independently demonstrate
the necessary photon-number balance.

An initially compatible Planck population does not,
however, require continuous photon creation merely
to preserve its form during subsequent evolution.

\subsection*{Falsification Criteria}

\begin{itemize}

\item \textbf{The Mechanical Reality of the Thermal Receiver:}

The proposed receiving modes must possess a physically
derived dispersion, state-counting measure, and energy
normalization. A two-dimensional interface cannot be
assigned three-dimensional thermal multiplicity without
an independently established interface density or
equivalent embedding.

Likewise, a gapped receiver must be evaluated using
its actual excitation spectrum, rather than a gapless
approximation outside its domain of validity.

\textit{Falsification:} A receiving construction that
violates these requirements, or cannot conserve energy,
momentum, and recoil through its proposed interactions,
fails as the specified microscopic thermal mechanism.

The carrier-projected irreversible response and the
thermal density of states are distinct requirements.
A receiver proposed for the local Ohmic photon-relaxation
channel must also produce the corresponding irreversible
spectrum; other thermal mechanisms need not share that
particular infrared law.

\item \textbf{The CMB Spectrum and Absolute Temperature:}

The completed Fossil-CMB mechanism must reproduce the
observed blackbody spectrum, absolute temperature,
and permitted spectral distortions from independently
specified initial conditions and constitutive properties.

\textit{Falsification:} If the resulting spectrum or
temperature falls outside observational constraints
after the microscopic parameters and thermal history
have been fixed independently, the proposed thermal
realization is falsified.

The historical near-$2.726\,\mathrm K$ conditional
evaluation is not an independent TFT prediction.
The observed temperature cannot be used to calibrate
an unknown receiver parameter and subsequently
presented as confirmation of that same calculation.

\item \textbf{Thermalization and the Survival of Cosmic Light:}

The early thermal interaction must be sufficiently
effective to establish the required radiation state,
yet the later universe must remain compatible with
the observed transparency, spectral precision,
directional structure, and polarization of surviving
astronomical photons.

\textit{Falsification:} If the calculated thermal
rates cannot establish the required spectrum during
the assigned epoch, or if subsequent interactions
produce excessive opacity, scattering, spectral
diffusion, polarization changes, or secondary radiation,
the proposed thermal history is contradicted.

A fossil radiation population need not remain in
active thermal equilibrium indefinitely. The
essential requirement is that equilibration occurs
when the proposed cosmic history requires it.

\item \textbf{The Temperature History Beyond the Observer:}

The Fossil Route also assigns a physical evolution
to the radiation temperature. Its predictions must
therefore remain compatible with temperature
diagnostics obtained from remote cosmic environments.

\textit{Falsification:} A completed thermal and
observer model that cannot reproduce independently
analyzed remote temperature measurements within
their uncertainties fails its predicted thermal
history.

Such comparisons must include the excitation,
collision, and radiative-transfer physics used
to infer the remote temperature. An inferred
temperature is not automatically a direct
measurement of the entire cosmic past.

\end{itemize}

\subsection*{The Boundary of Thermal Completion}

The central challenge is not to reproduce a numerical
temperature through a convenient energy partition.
It is to derive the receiving degrees of freedom,
their primordial loading, the photon interaction,
and the subsequent thermal evolution from the
underlying TFT mechanics.

The ideal phase-locked Hard Floor cannot simply be
assigned an unaccounted finite-inertia thermal bath.
The receiving dynamics and their energy ledger must
belong to a physically specified interaction.

Failure of a particular receiver would reject that
thermal construction, not independently established
FCC propagation or gravitational mechanics. Likewise,
the absence of a completed receiving mechanism is
an unresolved constitutive problem, not evidence
that every possible thermal realization has failed.

\medskip
\noindent\textbf{The fossil radiation must preserve
the history of its mechanical origin; its observed
spectrum and temperature must provide the test,
never the calibration used to manufacture the result.}

\section{Present Completion Status}

The present TFT CMB construction has the following status:

\begin{enumerate}

    \item
    \textbf{Coherent photon transport:}
    inherited from
    Appendices~\ref{app:photon_dispersion_relaxation}
    and~\ref{app:vacuumdamping_cosmicredshift}.
    A nonzero constant
    \(\Gamma_{\mathrm D}\)
    requires an independently established carrier-projected irreversible
    continuum.

    \item
    \textbf{Quantum-consistent energy transfer:}
    complete quantized receiving events can reproduce a smooth
    coarse-grained drift through their first statistical moment.
    The complete microscopic transition kernel remains to be specified.

    \item
    \textbf{Thermodynamic state counting:}
    for independently established three-dimensional gapless linear branches,
    \[
        \mathcal D/V\propto\Omega^2,
        \qquad
        u\propto T^4,
        \qquad
        C_V\propto T^3.
    \]

    \item
    \textbf{Dissipative projection:}
    the thermodynamic density of states and the carrier-projected dissipative
    spectrum are distinct physical objects.

    \item
    \textbf{Planck spectrum:}
    a zero-chemical-potential Planck state follows conditionally from a
    genuine photon-number-changing interaction satisfying detailed balance.

    \item
    \textbf{Fossil initial state:}
    catastrophic decompression provides the proposed origin of the primordial
    radiation, while
    \(T_{\mathrm a}\sim3000~{\rm K}\)
    is retained as the post-annealing thermal photon boundary.

    \item
    \textbf{Closed Fossil equilibrium:}
    Equation~\eqref{eq:AT_fossil_temperature_general}
    gives the common equilibrium temperature once the initial photon state,
    initial receiver state, and receiver thermodynamics are independently
    specified.

    \item
    \textbf{Finite active-to-latent interface:}
    the conditional assignment
    \(F_{\mathrm I}=1\)
    together with the proposed stiffness inheritance gives
    \[
        v_{\mathrm I}
        \simeq
        20.95~{\rm km\,s^{-1}}.
    \]
    If the same receiver is a distinct three-dimensional gapless branch with
    sufficiently small initial loading, the corresponding unrounded TFT
    quantities give
    \[
        T_{\mathrm f}
        \simeq
        2.726~{\rm K}.
    \]
    This provides an illustrative conditional consistency test of the
    finite-interface realization.

    \item
    \textbf{Interface infrared structure:}
    \(r_{\mathrm{QG}}\)
    can normalize the candidate interface inertia but cannot by itself
    generate the required low-frequency scale.

    \item
    \textbf{Internal phase branch:}
    phase-difference symmetry provides a conditional route to a gapless
    three-dimensional collective branch if finite phase inertia exists, but
    the regular local derivative coupling examined here is super-Ohmic and
    does not close the constant-\(\Gamma_{\mathrm D}\) problem.

    \item
    \textbf{Primordial receiver loading:}
    remains to be obtained from the decompression and annealing history.

    \item
    \textbf{Open thermal sink:}
    State-II structures provide a candidate capture architecture, while their
    microscopic absorption law and
    \(\Gamma_{\mathrm{II}}^{(\mathrm{cap})}\)
    remain to be derived.

    \item
    \textbf{Driven thermal fixed point:}
    \(T_{\mathrm{fp}}\)
    is conditionally determined once the thermal source, state-counting
    normalization, and sink law are independently specified.

    \item
    \textbf{Absolute CMB temperature:}
    the general energy-conserving Fossil closure is established conditionally,
    while the remaining microscopic receiver, loading, and thermalization
    structure determines whether a unique absolute temperature follows.

    \item
    \textbf{Astronomical observer map:}
    the relation between the local thermal state and
    \(T_{\mathrm{meas}}\)
    remains to be derived.

\end{enumerate}

The compact status is
\begin{equation}
    \begin{aligned}
        &\text{Fossil-radiation origin: proposed},
        \\
        &\text{quantized receiving architecture: mechanically admissible},
        \\
        &\text{thermal state counting: conditional},
        \\
        &\text{closed Fossil equilibrium: conditionally derived},
        \\
        &\text{finite-interface }2.726~{\rm K}
        \text{ endpoint: illustrative conditional specialization},
        \\
        &\text{Planck equilibrium: conditional on number-changing
        detailed balance},
        \\
        &\text{microscopic irreversible coupling: open},
        \\
        &\text{interface gap and primordial loading: open},
        \\
        &\text{absolute CMB temperature: not presently derived}.
    \end{aligned}
    \label{eq:AT_status_summary}
\end{equation}


\section{Summary}

\hyperref[app:cmb_thermal_completion]{Appendix~\ref*{app:cmb_thermal_completion}} establishes that three microscopic objects relevant to CMB
completion must remain distinct:
\begin{equation}
    \text{thermodynamic density of states}
    \;\neq\;
    \text{carrier-projected dissipative spectrum}
    \;\neq\;
    \text{thermal transfer kernel}.
    \label{eq:AT_summary_three_objects}
\end{equation}

For specified three-dimensional gapless linear receiving branches,
\begin{equation}
    \frac{
        \mathcal D_{\mathrm{AV}}(\Omega)
    }{
        V
    }
    =
    \frac{
        \Omega^2
    }{
        2\pi^2
    }
    \sum_j
    v_j^{-3},
\end{equation}
and
\begin{equation}
    u_{\mathrm{rec}}(T)
    =
    \frac{
        \pi^2k_{\mathrm B}^4
    }{
        30\hbar^3
    }
    \left(
        \sum_j
        v_j^{-3}
    \right)
    T^4.
    \label{eq:AT_summary_receiver_thermodynamics}
\end{equation}

For a closed Fossil photon--receiver subsystem,
\begin{equation}
    \boxed{
    T_{\mathrm{eq}}^4
    =
    \frac{
        2c_{\mathrm E}^{-3}
        T_{\gamma,0}^4
        +
        \left(
            \displaystyle\sum_j
            v_j^{-3}
        \right)
        T_{\mathrm{rec},0}^4
    }{
        2c_{\mathrm E}^{-3}
        +
        \displaystyle\sum_j
        v_j^{-3}
    }.
    }
    \label{eq:AT_summary_fossil_temperature}
\end{equation}

This relation is the central energy-conserving Fossil endpoint for
independently specified initial states and receiving thermodynamics.

The finite active-to-latent construction provides a particular conditional
specialization.

For
\begin{equation}
    F_{\mathrm I}=1,
\end{equation}
the proposed topology and stiffness inheritance give
\begin{equation}
    v_{\mathrm I}
    \simeq
    20.95~{\rm km\,s^{-1}},
\end{equation}
and the corresponding scalar thermodynamic weighting is
\begin{equation}
    \frac{
        a_{\mathrm I}
    }{
        a_\gamma
    }
    \simeq
    1.47\times10^{12}.
\end{equation}

If the interface receiver is a distinct three-dimensional gapless branch with
sufficiently small primordial loading, then
\begin{equation}
    T_{\mathrm f}
    =
    T_{\mathrm a}
    \left(
        1+
        \frac{
            a_{\mathrm I}
        }{
            a_\gamma
        }
    \right)^{-1/4}
\end{equation}
gives, using the unrounded TFT microscopic quantities and
\(T_{\mathrm a}=3000~{\rm K}\),
\begin{equation}
    T_{\mathrm f}
    \simeq
    2.726~{\rm K}.
\end{equation}

The measured present CMB temperature is not used to determine
\(F_{\mathrm I}\),
\(v_{\mathrm I}\),
or the thermal coefficient in this calculation.

This near-\(2.726~{\rm K}\) value is therefore retained as an illustrative
conditional consistency result of the finite-interface specialization.

The convergence radius
\(r_{\mathrm{QG}}\)
enters this construction as a microscopic density and inertia normalization,
not as the direct source of the infrared thermal frequency.

A Planck electromagnetic state additionally requires
\begin{equation}
    \frac{
        \Gamma_{\mathrm{em}}
    }{
        \Gamma_{\mathrm{abs}}
    }
    =
    \exp
    \left(
        -
        \frac{
            \hbar\omega
        }{
            k_{\mathrm B}T
        }
    \right),
    \label{eq:AT_summary_detailed_balance}
\end{equation}
together with photon-number-changing dynamics.

For an open thermal subsystem,
\begin{equation}
    C_V^{(\mathrm{th})}(T)\dot T
    =
    \mathcal P_{\mathrm{in}}
    -
    \mathcal P_{\mathrm{sink}}.
    \label{eq:AT_summary_open_balance}
\end{equation}

The resulting driven fixed-point temperature
\(T_{\mathrm{fp}}\)
is physically distinct from the closed Fossil equilibrium temperature
\(T_{\mathrm{eq}}\).

The Hard Floor retains its role as the zero-inertia constraint,
phase-reference, and information-memory sector. Finite thermal state
counting, heat capacity, and irreversible dissipation reside in finite-energy
Active-Vacuum, active-interface, radiative, or structural degrees of freedom.

The resulting hierarchy is
\begin{equation}
    \begin{aligned}
        \text{catastrophic decompression}
        &\longrightarrow
        \text{post-annealing fossil radiation},
        \\
        \text{coherent photon transport}
        &\longrightarrow
        \text{finite-energy receiving modes},
        \\
        \text{receiving modes}
        &\rightleftharpoons
        \text{thermal photon field},
        \\
        \text{thermalized sector}
        &\longrightarrow
        \text{macroscopic thermal balance}.
    \end{aligned}
    \label{eq:AT_summary_hierarchy}
\end{equation}

\hyperref[app:cmb_thermal_completion]{Appendix~\ref*{app:cmb_thermal_completion}} therefore supplies an energy-conserving Fossil-CMB framework and
a specific conditional finite-interface specialization without using the
measured CMB temperature as an upstream calibration.

The remaining microscopic bridges are the identity and three-dimensional
support of the receiving interface field, its stiffness and infrared gap, its
primordial loading, the photon--receiver thermalization interaction, the
carrier-projected irreversible continuum required for nonzero
\(\Gamma_{\mathrm D}\), and the final matter--radiation observer map.

The general energy-conserving two-sector endpoint is the principal
thermodynamic result retained from the Fossil-CMB analysis. The
near-\(2.726~{\rm K}\) interface endpoint is retained as a specific conditional
specialization to be revisited when the corresponding microscopic
Active-Vacuum interface dynamics are developed.

The observed CMB spectrum, absolute temperature, spectral-distortion limits,
and temperature measurements at nonzero astronomical redshift remain
downstream empirical tests of the completed TFT thermal sector rather than
inputs used to determine its unresolved microscopic parameters.